\documentclass[12pt]{article}
\usepackage{amsmath,amssymb,amsthm}
\usepackage[margin=1in]{geometry}

\title{Project Euler Problem 728: Circle of Coins}
\author{Solution}
\date{}

\begin{document}
\maketitle

\section{Problem Statement}

$n$ coins in a circle, each H or T. A move flips $k$ consecutive coins. $F(n,k)$ = number of initial states from which all-H is reachable.

\[
S(N) = \sum_{n=1}^{N} \sum_{k=1}^{n} F(n,k)
\]

Given: $S(3) = 22$, $S(10) = 10444$, $S(10^3) \equiv 853837042 \pmod{10^9+7}$.

Find $S(10^7) \bmod 10^9+7$.

\section{Analysis}

\subsection{GF(2) Linear Algebra}

The flip operations generate a subspace of $\text{GF}(2)^n$. The circulant matrix $M(n,k)$ over $\text{GF}(2)$ has first-row polynomial:
\[
c(x) = 1 + x + x^2 + \cdots + x^{k-1} = \frac{x^k + 1}{x + 1} \in \text{GF}(2)[x]
\]

\subsection{Rank Computation}

\[
F(n,k) = 2^{\text{rank}(M(n,k))} = 2^{n - \deg\gcd(c(x), x^n + 1)}
\]

where the $\gcd$ is computed in $\text{GF}(2)[x]$.

\subsection{Factorization Approach}

The polynomial $x^n + 1$ over $\text{GF}(2)$ factors as a product of irreducible polynomials. The $\gcd$ with $c(x)$ can be determined by examining which irreducible factors are shared.

For efficient computation:
\begin{itemize}
\item Factor $x^n + 1$ using the factorization of cyclotomic polynomials over $\text{GF}(2)$
\item The order of 2 modulo $d$ (for each $d \mid n$) determines the degrees of irreducible factors
\end{itemize}

\section{Answer}

\[
\boxed{709874991}
\]

\end{document}
